% TEMPLATE-MILC-v3.tex - plantilla maestra UNICA de las guias MILC v3.
%
% La usan las cuatro familias de guias, cada una con su driver:
%   · Grados 6-11          -> scripts/build-guias-g11.py
%   · Bebras               -> scripts/build-guias-bebras.py
%   · Semillero            -> scripts/build-guias-semillero.py
%   · Territorio interior  -> scripts/build-guias-territorio-interior.py
%
% Antes habia tres copias casi identicas (~400 lineas iguales, 6 distintas):
% cada mejora habia que aplicarla tres veces, y en la practica no se hacia
% (el motor de recursos vivia solo en dos de ellas). Lo que varia entre
% familias esta parametrizado en 6 placeholders que cada driver llena
% (se nombran aqui SIN su sintaxis real: el builder reemplaza por texto plano,
% tambien dentro de los comentarios, y un valor multilinea rompe el preambulo):
%
%   HEADER_IZQ        encabezado izquierdo de pagina
%   HEADER_DER        encabezado derecho de pagina
%   PORTADA_ETIQUETA  rotulo sobre el numero grande (GRADO/MOMENTO/MODULO)
%   PORTADA_SERIE     linea de serie de la portada
%   PORTADA_ID        identificador de la guia en la portada
%   PORTADA_CIRCULO   el circulito de grado (°), o vacio si no aplica
%   PDF_TITULO        titulo en texto plano (sin \textbf ni comillas TeX) para
%                     los metadatos del PDF (/Title); lo produce lib_xelatex.texto_pdf
%
% Composicion (auditoria 2026-09-03): las bandas de estacion piden espacio
% (\Needspace*) y prohiben el salto justo despues (\nopagebreak), asi no quedan
% huerfanas al pie; las cajas largas (softbox, drawbox, infoband, puentebox) son
% `breakable`, asi una caja de media pagina ya no arrastra todo a la siguiente
% dejando la anterior al 40 %. Todo texto sobre mostaza o verde va en milcNegro
% (blanco sobre #E5B400 da 1,93:1 y sobre #58A923 2,95:1; ninguno pasa WCAG AA).
% El resto de claves las documenta content/guias/_SCHEMA.md. El script valida
% que no queden placeholders sin reemplazar antes de escribir el archivo final.

% Etiquetado PDF (latex-lab): /Lang, estructura y texto alternativo de las
% figuras, para lectores de pantalla. Debe ir ANTES de \documentclass.
\DocumentMetadata{lang=es-CO,tagging=on}
\documentclass[11pt,letterpaper]{article}
% headheight=24pt: la cabecera derecha lleva dos lineas (identidad de la guia
% y estacion actual). headsep se acorta en la misma medida para que el bloque
% de texto quede exactamente donde estaba con headheight=12pt.
\usepackage[margin=2.0cm,headheight=24pt,headsep=13pt]{geometry}
\usepackage[table]{xcolor}
\usepackage{fontspec}
\usepackage[spanish]{babel}
\usepackage{tikz}
% Librerías TikZ para los diagramas de las guías (bloque `recursos:`):
%  · babel  → babel en español vuelve activos caracteres como «>», y TikZ los usa
%             en sus opciones (>=stealth, ->). Sin esta librería, cualquier
%             diagrama con flechas revienta con «\language@active@arg> has an
%             extra }». Debe ir ANTES de que se dibuje nada.
%  · positioning        → sintaxis «below left=2pt and 3pt».
%  · decorations.pathreplacing → llaves ({brace}) para señalar rangos.
\usetikzlibrary{babel,positioning,decorations.pathreplacing}
\usepackage{graphicx}
% Circuitos: el trabajo de electrónica se hace hoy con micro:bit y sus sensores
% integrados (MakeCode, sin cableado), así que no se necesitan esquemáticos.
% Si algún día entran montajes con protoboard, basta descomentar esta línea:
% los diagramas .tex/.tikz declarados en `recursos:` ya se incrustan solos.
% \usepackage{circuitikz}
\usepackage[most]{tcolorbox}
\usepackage{tabularx}
\usepackage{array}
\usepackage{fancyhdr}
\usepackage{titlesec}
\usepackage[document]{ragged2e}  % bandera a la derecha en todo el cuerpo
\usepackage{enumitem}
\usepackage{microtype}
\usepackage{etoolbox}
\usepackage{needspace}
% hyperref al final del preambulo: metadatos del PDF (titulo, autor, idioma,
% familia), marcadores por estacion (\pdfbookmark) y enlaces sin recuadro.
\usepackage[hidelinks,
  pdftitle={Duelo: no hay etapas que cumplir, hay compañía que organizar},
  pdfauthor={PhD. Álvaro Cárdenas Orozco},
  pdfsubject={Territorio interior · Educación socioemocional · Grado 10° · Momento 5 · MILC v3},
  pdflang={es-CO},
  bookmarksopen=true,bookmarksnumbered=false]{hyperref}

% Helvetica Neue no trae la flecha U+2192, asi que cada «→» escrito literalmente
% en el YAML se compilaba a NADA, en silencio (462 flechas en 61 guias: rutas del
% tipo «paleta → tipografia → lockup» salian rotas). Mapeamos el caracter a su
% version matematica, que si tiene glifo. Asi el autor puede seguir escribiendo
% «→» comodamente en el YAML.
\usepackage{newunicodechar}
\newunicodechar{→}{\ensuremath{\rightarrow}}
% Mismo problema con los simbolos de marca y vinetas: el «✓» de «una palomita ✓»
% salia como cuadrito vacio en el PDF de todas las guias que lo usaban.
\usepackage{pifont}
\newunicodechar{✓}{\ding{51}}
\newunicodechar{✗}{\ding{55}}
\newunicodechar{✔}{\ding{52}}
\newunicodechar{•}{\textbullet}
\newunicodechar{≥}{\ensuremath{\geq}}
\newunicodechar{≤}{\ensuremath{\leq}}

\setmainfont{Helvetica Neue}
\setsansfont{Helvetica Neue}
\setlength{\parindent}{0pt}
\setlength{\parskip}{6pt}
% Sin particion de palabras: en 6.º-7.º una palabra cortada por guion al final
% de linea («Comuni-cacion») frena la lectura mucho mas que una linea corta.
\hyphenpenalty=10000
\exhyphenpenalty=10000
\pretolerance=2000
\tolerance=3000
\setlength{\emergencystretch}{3em}
\linespread{1.10}
\renewcommand{\arraystretch}{1.25}
\setlist{leftmargin=5mm,itemsep=1mm,topsep=1mm}
\newcolumntype{Y}{>{\RaggedRight\arraybackslash}X}

\definecolor{milcMagenta}{HTML}{E6007E}
\definecolor{milcVino}{HTML}{5A0038}
\definecolor{milcTurquesa}{HTML}{0093A5}
\definecolor{milcVerde}{HTML}{58A923}
\definecolor{milcMostaza}{HTML}{E5B400}
\definecolor{milcOcre}{HTML}{B86600}
\definecolor{milcNegro}{HTML}{241816}
\definecolor{milcGris}{HTML}{F7F7F4}
\definecolor{milcLinea}{HTML}{E8E2DD}
\definecolor{milcGradePrimary}{HTML}{0D9488}
\definecolor{milcGradeDark}{HTML}{115E59}
\definecolor{milcGradeSoft}{HTML}{CCFBF1}
\definecolor{milcGradeAccent}{HTML}{CFE7FF}
\definecolor{milcGradeText}{HTML}{FFFFFF}
\color{milcNegro}

\pagestyle{fancy}
\fancyhf{}
% Cabecera de dos lineas: identidad de la guia arriba y, a la derecha y en
% gris, la estacion en curso (\markright desde \milcestacion). Antes de la
% primera estacion la segunda linea va vacia; el \strut mantiene la altura
% para que la linea de arriba no baile entre paginas.
\lhead{\small Territorio interior · Educación socioemocional\\[-1pt]{\footnotesize\strut}}
\rhead{\small Grado 10° · Momento 5 · MILC v3\\[-1pt]{\footnotesize\strut\color{milcNegro!65}\rightmark}}
\cfoot{\small\thepage}
\renewcommand{\headrulewidth}{0pt}

% Sobre mostaza (#E5B400) y verde (#58A923) el texto va en milcNegro; sobre el
% resto de la paleta, en blanco. Se usa dentro de fonttitle/fontupper, donde un
% \color posterior manda sobre coltitle.
\newcommand{\milctintasobre}[1]{%
  \ifboolexpr{test{\ifstrequal{#1}{milcMostaza}} or test{\ifstrequal{#1}{milcVerde}}}%
    {\color{milcNegro}}{\color{white}}}

\newtcolorbox{titlebox}[1]{
  enhanced,colback=#1,colframe=#1,arc=7mm,boxrule=0pt,
  left=7mm,right=7mm,top=3mm,bottom=3mm,
  before skip=8pt,after skip=8pt,
  fontupper=\sffamily\bfseries\Large\color{white}
}

\newtcolorbox{softbox}[3]{
  enhanced,breakable,pad at break=3mm,
  colback=#2,colframe=#1,borderline west={4pt}{0pt}{#1},
  arc=2mm,boxrule=.4pt,left=4mm,right=4mm,top=3mm,bottom=3mm,
  before skip=6pt,after skip=8pt,title={#3},coltitle=white,
  colbacktitle=#1,fonttitle=\sffamily\bfseries\milctintasobre{#1},fontupper=\RaggedRight,
  attach boxed title to top left={xshift=3mm,yshift=-2mm},
  boxed title style={arc=2mm, boxrule=0pt}
}

\newtcolorbox{drawbox}[2]{
  enhanced,breakable,pad at break=4mm,
  colback=white,colframe=#1,arc=4mm,boxrule=1pt,
  left=5mm,right=5mm,top=4mm,bottom=4mm,before skip=8pt,after skip=8pt,
  title={#2},coltitle=white,colbacktitle=#1,fonttitle=\sffamily\bfseries\milctintasobre{#1},
  fontupper=\RaggedRight,
  attach boxed title to top left={xshift=4mm,yshift=-2mm},
  boxed title style={arc=3mm, boxrule=0pt}
}

\newtcolorbox{infoband}[1]{
  enhanced,breakable,pad at break=2mm,colback=#1!8,colframe=#1!50,arc=2mm,boxrule=.5pt,
  left=4mm,right=4mm,top=2mm,bottom=2mm,before skip=4pt,after skip=4pt,
  fontupper=\small\RaggedRight
}

% Puente narrativo entre secciones (contrato editorial Conéctate).
% Da continuidad: "Acabas de X. Ahora vas a Y porque Z."
% Visualmente: banda izquierda en color del grado, texto en itálica pequeña.
\newtcolorbox{puentebox}{
  enhanced,breakable,pad at break=2mm,colback=milcGradeSoft,colframe=milcGradeDark,
  borderline west={3pt}{0pt}{milcGradeDark},
  arc=0mm,boxrule=0pt,left=5mm,right=4mm,top=2.5mm,bottom=2.5mm,
  before skip=4pt,after skip=8pt,
  fontupper=\itshape\small\color{milcNegro}\RaggedRight
}
% La flecha va en modo matematico a proposito: Helvetica Neue NO tiene el glifo
% U+2192, asi que \textrightarrow se compilaba a NADA («Missing character: There
% is no → in font Helvetica Neue») y el puente salia sin flecha. En modo
% matematico la flecha viene de la fuente matematica y si se dibuja.
\newcommand{\puente}[1]{%
  \begin{puentebox}%
    \textcolor{milcGradeDark}{$\rightarrow$}\hspace{2mm}#1%
  \end{puentebox}%
}

\newcommand{\linead}[1]{\noindent\color{milcLinea}\rule{#1}{0.5pt}\color{milcNegro}}
\newcommand{\respuesta}[1]{\par\vspace{2mm}\foreach \n in {1,...,#1}{\noindent\color{milcLinea}\rule{\linewidth}{0.5pt}\par\vspace{5mm}}\color{milcNegro}}
\newcommand{\checkbox}{\textcolor{milcGradeDark}{\fbox{\phantom{x}}}\hspace{5pt}}

% ─── Listas aireadas para las guías socioemocionales ────────────────────
% Componer las verificaciones como «\checkbox texto\\» deja el texto que salta
% de línea debajo del cuadrito y apelmaza el bloque. Estos entornos alinean la
% continuación y airean los ítems. Los usa Territorio Interior; cualquier
% familia puede hacerlo.
\newlist{checklista}{itemize}{1}
\setlist[checklista]{
  label=\textcolor{milcGradeDark}{\fbox{\phantom{x}}},
  leftmargin=9mm, labelsep=3.2mm, itemsep=2.4mm,
  topsep=2.5mm, parsep=0pt, partopsep=0pt, before=\RaggedRight
}
\newlist{pasoslista}{enumerate}{1}
\setlist[pasoslista]{
  label=\textcolor{milcGradeDark}{\sffamily\bfseries\arabic*.},
  leftmargin=8mm, labelsep=3mm, itemsep=2.8mm,
  topsep=1mm, parsep=0pt, partopsep=0pt, before=\RaggedRight
}

\newcommand{\phasecard}[4]{
\begin{tcolorbox}[enhanced,colback=#3,colframe=#2,arc=3mm,boxrule=.5pt,height=3.05cm,valign=center,halign=center,left=1.5mm,right=1.5mm,top=2mm,bottom=2mm]
{\sffamily\bfseries\fontsize{22}{24}\selectfont\textcolor{#2}{#1}}\par
{\sffamily\bfseries\small #4}
\end{tcolorbox}
}

% ── Piezas del rediseno 2026 (legibilidad 6.º-11.º) ───────────────────
% Banda de estacion: reemplaza el titlebox macizo. Titulo a la izquierda,
% dato practico (tiempo/modalidad) a la derecha, en una sola linea.
% Nunca huerfana: pide 9 lineas libres (o salta de pagina rellenando con
% \vfill) y prohibe el salto justo despues de la banda. Nueve y no seis:
% la caja partible que sigue necesita ~80pt (titulo adosado, relleno y dos
% lineas) para arrancar en la misma pagina; con menos, tcolorbox la manda
% entera a la siguiente y la banda se queda sola. El penalty 10000 va
% en `after`, ANTES del espacio posterior: un `after skip` (aunque sea 0pt)
% es un \vskip tras la caja, y ese glue es un punto de corte legal que un
% \nopagebreak escrito despues de \end{tcolorbox} ya no cierra. Cada banda
% deja marcador PDF y actualiza la cabecera derecha.
\newcounter{milcestacion}
\newcommand{\milcestacion}[2]{%
\Needspace*{9\baselineskip}%
\stepcounter{milcestacion}%
\pdfbookmark[1]{#2}{milcestacion\themilcestacion}%
\markright{#2}%
\begin{tcolorbox}[enhanced,colback=#1,colframe=#1,arc=2mm,boxrule=0pt,
  left=5mm,right=5mm,top=2.6mm,bottom=2.6mm,before skip=11pt,
  after={\par\nopagebreak\vskip7pt}]
{\sffamily\bfseries\fontsize{15}{18}\selectfont\milctintasobre{#1}#2}
\end{tcolorbox}}

% Tarjeta de paso numerada: sustituye las tablas «Pilar / Aplicacion», que
% eran formato de consulta docente, no de estudiante. El numero va en un
% circulo sobre el borde para que la secuencia se lea de un vistazo.
\newcommand{\milcpaso}[3]{%
\Needspace{4\baselineskip}%
\begin{tcolorbox}[enhanced,colback=white,colframe=milcLinea,arc=1.5mm,boxrule=.9pt,
  left=14mm,right=4mm,top=3mm,bottom=3mm,before skip=4pt,after skip=4pt,
  fontupper=\RaggedRight,
  overlay={\node[anchor=north west,circle,fill=#2,text=white,inner sep=0pt,
    minimum size=7.4mm,font=\sffamily\bfseries\small]
    at ([xshift=3.6mm,yshift=-3.2mm]frame.north west) {#1};}]
#3
\end{tcolorbox}}

% Casilla marcable de tamano util con lapiz (la anterior era \fbox{\phantom{x}},
% de alto variable segun el texto que la seguia).
\newcommand{\milcmarca}{\framebox[3.8mm]{\rule{0pt}{3.4mm}}}

% Fila marcable de lista suelta (checks de la estacion 1).
\newcommand{\milccheckrow}[1]{%
\par\vspace{1.8mm}\noindent
\makebox[6.4mm][l]{\raisebox{-.55mm}{\textcolor{milcNegro}{\milcmarca}}}%
\parbox[t]{\dimexpr\linewidth-6.4mm\relax}{\RaggedRight #1}\par}

% Celda marcable dentro de la rubrica (tabla de autoevaluacion). Misma
% sangria francesa, pero pensada para vivir dentro de una columna X.
\newcommand{\milcopcion}[1]{%
\makebox[5.2mm][l]{\raisebox{-.55mm}{\textcolor{milcNegro}{\milcmarca}}}%
\parbox[t]{\dimexpr\linewidth-5.2mm\relax}{\RaggedRight #1}}

% ── Recursos visuales de la guía ──────────────────────────────────────
% Declarados en el bloque `recursos:` del YAML y emitidos por el builder.
% \guiaFigura   → imagen rasterizada o PDF (foto, carta celeste, montaje).
% \guiaDiagrama → fuente TikZ/circuitikz (.tex) incrustada como vector:
%                 es la vía para circuitos y esquemáticos editables.
% [#1] = texto alternativo (alt) · #2 = ruta relativa al .tex · #3 = pie
% El alt llega al PDF etiquetado (/Alt) y lo lee el lector de pantalla.
\newcommand{\guiaFigura}[3][]{%
\begin{center}
\includegraphics[width=0.88\linewidth,height=0.38\textheight,keepaspectratio,alt={#1}]{#2}\\[1.5mm]
{\footnotesize\itshape\textcolor{milcNegro!75}{#3}}
\end{center}
}

\newcommand{\guiaDiagrama}[2]{%
\begin{center}
\input{#1}\\[1.5mm]
{\footnotesize\itshape\textcolor{milcNegro!75}{#2}}
\end{center}
}

\begin{document}
\thispagestyle{empty}

% ── Portada ───────────────────────────────────────────────────────────
% Antes: un rectangulo de color a pagina completa. Se veia bien en pantalla,
% pero en la fotocopiadora del colegio se comia el toner y salia gris sucio.
% Ahora la banda ocupa el tercio superior y el resto va en blanco.
\begin{tikzpicture}[remember picture,overlay]
  \path[fill=milcGradePrimary]
    (current page.north west) rectangle ([yshift=-10.6cm]current page.north east);

  \node[anchor=north west,text=milcGradeText] at ([xshift=2.0cm,yshift=-1.55cm]current page.north west)
    {\sffamily\bfseries\fontsize{12}{14}\selectfont MOMENTO};
  \node[anchor=north west,text=milcGradeText,opacity=.85] at ([xshift=2.0cm,yshift=-2.35cm]current page.north west)
    {\sffamily\bfseries\fontsize{8.5}{10}\selectfont TERRITORIO INTERIOR · SOCIOEMOCIONAL};
  \node[anchor=north east,text=milcGradeText,opacity=.85,align=right] at ([xshift=-2.0cm,yshift=-1.55cm]current page.north east)
    {\sffamily\bfseries\fontsize{9}{12}\selectfont I.E. Sor María Juliana\\Cartago, Valle del Cauca};

  \node[anchor=west,text=milcGradeText] (gradonum) at ([xshift=1.85cm,yshift=-7.1cm]current page.north west)
    {\sffamily\bfseries\fontsize{112}{112}\selectfont 5};
  % El circulito de grado (°) solo aplica a las guias de grado («11°»); en
  % Bebras y Semillero el numero grande es un momento/modulo, no un grado.
  % Se posiciona relativo al nodo `gradonum`, no en coordenadas absolutas.
  

  \node[anchor=west,text=milcGradeText,text width=11.4cm,align=left] at ([xshift=1.5cm,yshift=0.5cm]gradonum.east)
    {
     {\sffamily\bfseries\fontsize{9}{11}\selectfont\MakeUppercase{Territorio interior · Socioemocional}}\\[3mm]
     {\sffamily\bfseries\fontsize{21}{25}\selectfont\setlength{\baselineskip}{27pt}Duelo\\[4pt]no hay etapas\\[4pt]que cumplir}\\[3.5mm]
     {\sffamily\bfseries\fontsize{15}{18}\selectfont Grado 10° · Momento 5}
    };
\end{tikzpicture}

\vspace*{9.4cm}
\vfill

{\sffamily\bfseries\fontsize{8}{10}\selectfont\textcolor{milcGradeDark}{LO QUE TE LLEVAS HOY}}

\begin{tcolorbox}[enhanced,colback=white,colframe=milcNegro,arc=2mm,boxrule=1.6pt,
  left=5mm,right=5mm,top=4mm,bottom=4mm,before skip=3pt,after skip=12pt,
  fontupper=\RaggedRight\fontsize{11}{15.5}\selectfont]
Tu calendario de acompañamiento: una pérdida —propia o de alguien cercano— con lo que de verdad ayudó y lo que no, la lista de frases que hacen daño con su reemplazo, y un plan de compañía concreto con días marcados, incluido el que todos olvidan.
\end{tcolorbox}

\begin{tcolorbox}[enhanced,colback=milcGris,colframe=milcGris,arc=2mm,boxrule=0pt,
  left=5mm,right=5mm,top=3.5mm,bottom=3.5mm,after skip=12pt]
{\sffamily\bfseries\fontsize{8}{10}\selectfont\textcolor{milcNegro!55}{ESTA GUÍA ES DE}}\\[3mm]
\begin{tabularx}{\linewidth}{X p{3.2cm} p{3.2cm}}
\linead{\linewidth} & \linead{3.0cm} & \linead{3.0cm}\\[-1mm]
{\fontsize{8.5}{10}\selectfont\textcolor{milcNegro!55}{Nombre completo}} &
{\fontsize{8.5}{10}\selectfont\textcolor{milcNegro!55}{Grupo}} &
{\fontsize{8.5}{10}\selectfont\textcolor{milcNegro!55}{Fecha}}\\
\end{tabularx}
\end{tcolorbox}

\vfill

\noindent\textcolor{milcLinea}{\rule{\linewidth}{1pt}}\\[2mm]
{\sffamily\bfseries\fontsize{9.5}{12}\selectfont Autor: Dr. Álvaro Cárdenas Orozco}\\[1mm]
{\sffamily\fontsize{8.5}{11}\selectfont\textcolor{milcNegro!55}{Metodología MILC · Apertura ancestral · Duelo · Triángulo Dussel-Estoicismo-Floridi}}

\newpage

\section*{Apertura: Duelo: no hay etapas que cumplir, hay compañía que organizar}

\begin{softbox}{milcGradeDark}{milcGradeSoft}{Saber ancestral}
En el Pacífico colombiano, cuando alguien muere, \textbf{el duelo tiene una forma, un ritmo y un final marcado} ---y nada de eso queda librado a que a alguien se le ocurra pasar por la casa---. Después del entierro empieza el \textbf{novenario}: \textbf{durante nueve días se reza el rosario, y la comunidad se reúne \emph{al caer la tarde}}. \emph{Nueve días. Todos. A la misma hora.} \textbf{Y al noveno llega el \textbf{levantamiento de tumba}, que es la despedida definitiva}: se desmonta lo que se había montado y el ritual se cierra. \emph{Los gualíes ---el ritual propio de cuando muere un niño---, los alabaos y los levantamientos de tumba fueron declarados \textbf{Patrimonio Inmaterial de la Nación en 2014}.} \textbf{Ahora fíjate en lo que ese sistema resuelve, porque no es lo que parece.} \emph{No organiza \textbf{el sentimiento} ---nadie le dice a la viuda cómo debe sentirse el martes---.} \textbf{Organiza \emph{la compañía}:} \emph{deja establecido quién viene, cuándo viene y hasta cuándo}. \textbf{Y ahí está la diferencia con lo que solemos hacer:} \emph{nosotros dejamos el sentimiento con horario ---«ya deberías estar mejor»--- y la compañía sin horario ---«cualquier cosa me avisas»---.} \textbf{Está exactamente al revés.}
\end{softbox}

\begin{softbox}{milcTurquesa}{milcGris}{Saber contemporáneo}
Casi todo el mundo ha oído que el duelo tiene \textbf{cinco etapas} ---negación, ira, negociación, depresión y aceptación--- y que hay que pasarlas en orden hasta llegar a la última. \textbf{Eso no es lo que dice la investigación, y ni siquiera es lo que dijo quien las propuso.} \emph{Elisabeth Kübler-Ross describió cinco estados como referencia para \textbf{entender} lo que le pasa a alguien, no como una escalera que haya que subir en orden;} \textbf{ella misma reconoció que su modelo fue malinterpretado durante décadas.} \textbf{Y cuando se ha puesto a prueba, no aguanta.} \emph{Un estudio de la Universidad de Yale entre 2000 y 2003 encontró que \textbf{la aceptación era la respuesta más común en todos los momentos} después de la pérdida ---no solo al final---, lo que contradice la idea de que sea «la última etapa».} \emph{Un contraste empírico del modelo secuencial mostró un \textbf{ajuste pobre} a los datos, mientras que un modelo no secuencial sí ajustaba bien.} \textbf{Y George Bonanno, tras dos décadas de estudios, concluyó que no hay evidencia que sostenga la teoría de las etapas.} \emph{La consecuencia práctica es enorme:} \textbf{nadie está «atrasado» en su duelo, y a nadie hay que ayudarlo a pasar a la etapa siguiente.}
\end{softbox}

\begin{softbox}{milcMostaza}{milcGris}{Pregunta-puente}
\textit{El novenario organiza \textbf{quién viene y hasta cuándo}, y no le dice a nadie cómo debe sentirse. \textbf{Cuando alguien de tu entorno perdió algo importante, ¿en qué momento dejó de tener compañía} \ldots \textbf{y cuánto tiempo llevaba}?}
\end{softbox}

\begin{softbox}{milcVerde}{milcGris}{Saber y saber hacer}
Al terminar podrás: (1) \textbf{explicar} por qué las «cinco etapas» no son una ruta que haya que cumplir, con la evidencia; (2) \textbf{reconocer} qué frases hacen daño aunque se digan con cariño, y con qué reemplazarlas; (3) \textbf{entender} que lo que un duelo necesita no es apuro sino compañía organizada; (4) \textbf{planear} un acompañamiento concreto, con días marcados, incluido el momento en que todos los demás ya se fueron.
\end{softbox}



\small
\begin{tabularx}{\textwidth}{>{\bfseries}p{3.0cm}Y}
\rowcolor{milcGradePrimary!12}Autor & Dr. Álvaro Cárdenas Orozco\\
DBA & Comprende los fenómenos que afectan a su territorio, regula sus emociones ante ellos y acompaña a otros con empatía (transversal 6°--11°, articulado con la política «Escuela, territorio de vida» del MEN, Resolución 006519 de 2025, y la Ley 1620 de 2013)\\
Referentes & Primera ayuda psicológica (OMS, 2011/2012) · Directrices IASC (2007) · USGS y Servicio Geológico Colombiano · Pueblo nasa y la minga (CRIC) · Dussel · Estoicismo · Floridi\\
Producto final & Tu calendario de acompañamiento: una pérdida —propia o de alguien cercano— con lo que de verdad ayudó y lo que no, la lista de frases que hacen daño con su reemplazo, y un plan de compañía concreto con días marcados, incluido el que todos olvidan.\\
\end{tabularx}
\normalsize

\puente{Este momento sigue a los tres anteriores y es distinto: \textbf{aquí no hay nada que arreglar}. Se trata de aprender a estar al lado de alguien sin apurarlo.}

\milcestacion{milcGradeDark}{Ruta MILC: así vamos a aprender}
\begin{tabularx}{\textwidth}{>{\centering\arraybackslash}X>{\centering\arraybackslash}X>{\centering\arraybackslash}X>{\centering\arraybackslash}X}
\phasecard{1}{milcVerde}{milcGris}{Escucha\\[-1mm]\small Observo, escucho y pregunto.} &
\phasecard{2}{milcTurquesa}{milcGris}{Sistematización\\[-1mm]\small No lo apuro, lo acompaño.} &
\phasecard{3}{milcMagenta}{milcGris}{Praxis\\[-1mm]\small Construyo y aplico.} &
\phasecard{4}{milcVino}{milcGris}{Evaluación\\[-1mm]\small Reflexión liberadora.}
\end{tabularx}


\puente{Primero, una pérdida real ---tuya o cercana--- y la pregunta que casi nadie hace: \emph{¿qué ayudó de verdad?}}

\milcestacion{milcVerde}{Estación 1 de 3 · Escucha}

\begin{softbox}{milcVerde}{milcGris}{Escena para escuchar}
\textbf{Actividad 1 · IDENTIFICA --- Lo que de verdad ayudó.} \textbf{Nadie va a leer esta página.} \textbf{Primera parte.} Piensa en \textbf{una pérdida} ---tuya o de alguien cercano---. \emph{No tiene que ser una muerte:} \textbf{una separación en la casa, una mudanza, un amigo que se fue del país, una ruptura, un animal que se murió, algo que se perdió para siempre.} \textbf{Segunda parte.} Escribe \textbf{qué hizo la gente} y clasifícalo en dos columnas: \emph{lo que ayudó} y \emph{lo que no}. \textbf{Sé específico:} \emph{«me trajeron comida» es distinto de «me dijeron que fuera fuerte».} \textbf{Tercera parte, la que revela el patrón.} \emph{¿Cuánto tiempo duró el acompañamiento?} \textbf{Y después:} \emph{¿cuánto tiempo duró la tristeza?} \textbf{Compara los dos números.} \textbf{Y la última:} \emph{¿alguien te preguntó ---o le preguntó--- algo un mes después? ¿Y a los seis meses?}
\end{softbox}

{\sffamily\bfseries\fontsize{8}{10}\selectfont\textcolor{milcNegro!55}{MARCA CUANDO LO HAYAS HECHO}}
\milccheckrow{Escogiste una pérdida real, no necesariamente una muerte.}
\milccheckrow{Separaste \textbf{lo que ayudó} de \textbf{lo que no}, con hechos concretos.}
\milccheckrow{Comparaste cuánto duró la compañía con cuánto duró la tristeza.}

\begin{infoband}{milcVerde}


\textbf{En tu cuaderno (Actividad 1 · IDENTIFICA).} Título: «Actividad 1. Lo que de verdad ayudó». \textbf{Formato:} la pérdida + las dos columnas + los dos tiempos comparados + si alguien preguntó después. \textbf{Extensión:} media página. \textbf{Sabes que terminaste cuando} tienes \textbf{los dos números} escritos. Esta página es tuya: \textbf{nadie más tiene que leerla si tú no quieres}.
\end{infoband}

\puente{Ya la tienes. Ahora, por qué las etapas no son una ruta y qué es lo que un ritual como el novenario resuelve de verdad.}

\milcestacion{milcTurquesa}{Fase 2 · Sistematización: entender el duelo}

\begin{softbox}{milcTurquesa}{milcGris}{Cuatro claves sobre el duelo}
Cuatro claves para dejar de apurar a la gente y empezar a acompañarla.
\end{softbox}

\milcpaso{1}{milcTurquesa}{\textbf{Las cinco etapas no son una ruta, y quien las propuso no dijo que lo fueran.} \emph{Kübler-Ross describió cinco estados como referencia para entender, no como pasos en orden, y reconoció que se malinterpretó su modelo.} \textbf{Y la evidencia lo confirma:} \emph{el estudio de Yale (2000-2003) encontró que \textbf{la aceptación era la respuesta más común en todos los momentos}, no solo al final; un contraste empírico mostró que el modelo secuencial \textbf{ajusta mal} y el no secuencial ajusta bien; y Bonanno, tras dos décadas de investigación, no halló evidencia que sostenga la teoría.} \textbf{Consecuencia:} \emph{nadie está atrasado en su duelo, nadie está «estancado en la negación», y la frase «ya deberías estar en otra etapa» no tiene ningún respaldo.} \textbf{El duelo va y viene:} \emph{un día bueno no significa que ya pasó, y una recaída no significa retroceso.}}
\milcpaso{2}{milcTurquesa}{\textbf{No se pasa: se aprende a llevar.} \emph{La palabra «superar» hace daño porque promete que el otro va a desaparecer del todo, y eso no ocurre ---ni tendría por qué---.} \textbf{Lo que cambia no es el tamaño de la ausencia: es que la vida vuelve a crecer alrededor.} \emph{Al principio la pérdida ocupa todo; con el tiempo hay otras cosas ---y la pérdida sigue ahí, del mismo tamaño, pero ya no es lo único---.} \textbf{Por eso las fechas duelen años después y eso no es un síntoma de nada:} \emph{es que el vínculo no se borró}. \emph{Y esto también sirve para uno mismo:} \textbf{si a los seis meses todavía duele, no estás fallando.}}
\milcpaso{3}{milcTurquesa}{\textbf{El ritual no organiza el sentimiento: organiza la compañía.} \emph{Vuelve a mirar el novenario:} \textbf{nueve días, todos, a la misma hora, y un cierre marcado en el noveno.} \emph{Nadie le dice a nadie cómo debe sentirse ---eso queda libre---.} \textbf{Lo que queda establecido es \emph{quién viene y hasta cuándo}, que es justo lo que hoy queda al azar.} \emph{Y hay una segunda cosa que resuelve, y es enorme:} \textbf{le quita a la persona en duelo el trabajo de pedir}. \emph{Quien está mal casi nunca llama ---no tiene fuerza, o no quiere molestar---, así que un sistema donde \textbf{la gente simplemente llega} es infinitamente mejor que uno donde hay que avisar.}}
\milcpaso{4}{milcTurquesa}{\textbf{El día que todos olvidan.} \emph{En los primeros días sobra la gente:} \textbf{llega comida, llegan visitas, llega todo el mundo.} \emph{Y después, de golpe, no llega nadie ---justo cuando el ruido se acaba y empieza lo difícil---.} \textbf{Compara los dos números de tu Actividad 1 y vas a ver la brecha:} \emph{el acompañamiento dura días y la tristeza dura meses.} \textbf{Por eso el acompañamiento que de verdad sirve es el que aparece \emph{tarde}:} \emph{el mensaje del mes siguiente, la llamada a los seis meses, el «me acordé de él» en la fecha.} \emph{Y hay un detalle que casi nadie sabe y que agradece todo el que ha perdido a alguien:} \textbf{decir el nombre de la persona que murió.} \emph{Casi nadie lo dice, por miedo a hacer daño ---y el que quedó suele estar deseando oírlo---.}}

\begin{drawbox}{milcTurquesa}{Lo que se dice y con qué reemplazarlo}
\begin{pasoslista}
\item \textbf{«Ya está descansando» · «todo pasa por algo»} → \emph{«no sé qué decir, pero aquí estoy».}
\item \textbf{«Tienes que ser fuerte»} → \emph{«llora lo que necesites».}
\item \textbf{«Ya deberías estar mejor»} → \emph{«¿cómo va hoy?».}
\item \textbf{«Cualquier cosa me avisas»} → \emph{«paso el jueves a las cuatro».}
\item \textbf{Y la mejor de todas, que casi nadie usa:} \emph{decir el nombre de quien murió.}
\end{pasoslista}
\end{drawbox}

\begin{softbox}{milcMostaza}{milcGris}{Errores comunes y su corrección}
\textbf{Hablar de etapas:} nadie está atrasado ni estancado. \textbf{Decir «supéralo»:} no se pasa, se aprende a llevar. \textbf{Buscarle el lado bueno:} no le toca a nadie encontrarlo, y menos a los quince días. \textbf{«Cualquier cosa me avisas»:} el que está mal no llama. \textbf{Desaparecer a la semana:} ahí es donde empieza lo difícil. \textbf{Evitar el nombre:} el que quedó suele estar deseando oírlo. \textbf{Comparar duelos:} «yo perdí a mi abuelo y seguí» cierra la conversación.
\end{softbox}

\begin{infoband}{milcTurquesa}


\textbf{En tu cuaderno (Actividad 2 · EXPLICA).} Título: «Actividad 2. Las etapas no son una ruta». \textbf{Formato:} escribe qué encontraron el estudio de Yale y Bonanno, y explica por qué decir «ya deberías estar mejor» no tiene respaldo. \textbf{Extensión:} media página. \textbf{Sabes que terminaste cuando} puedes explicar la diferencia entre \textbf{organizar el sentimiento y organizar la compañía}.
\end{infoband}

\puente{Ya sabes que no se apura. Es hora de organizar la compañía, que es lo único que sí está en tus manos.}

\milcestacion{milcMagenta}{Fase 3 · Praxis: organizar la compañía}

\begin{softbox}{milcMagenta}{milcGris}{Cómo se acompaña un duelo sin apurarlo}
Aquí no se aprende qué decir. Se aprende a aparecer, que es lo que de verdad se recuerda.
\end{softbox}

\milcpaso{1}{milcMagenta}{\textbf{Lo que ayudó y lo que no (7 min).} Termina tus dos columnas y \textbf{fíjate en un patrón}: \emph{casi siempre lo que ayudó fueron \textbf{acciones} ---comida, compañía, alguien que se quedó, alguien que hizo un trámite--- y lo que no ayudó fueron \textbf{frases}}. \textbf{Ese solo hallazgo cambia la manera de acompañar a alguien.}}
\milcpaso{2}{milcMagenta}{\textbf{Las frases y su reemplazo (8 min).} Toma \textbf{cuatro} frases de tu columna de «no ayudó» y escribe con qué las reemplazarías. \textbf{Condición:} \emph{que el reemplazo no le pida nada a la persona en duelo} ---ni que responda, ni que agradezca, ni que se sienta de alguna manera---.}
\milcpaso{3}{milcMagenta}{\textbf{El calendario (12 min).} Escoge a \textbf{alguien real} que esté pasando por una pérdida ---o que la haya pasado este año--- y arma un \textbf{calendario de compañía} con \textbf{días concretos}: \emph{esta semana, dentro de un mes, y una fecha lejana} ---el cumpleaños, el aniversario, la fecha en que pasó---. \textbf{La regla que hace la diferencia:} \emph{nada de «cualquier cosa me avisas»; todo con día y hora.} \emph{Y no hace falta anunciarlo: se pone en el calendario del celular y ya.}}
\milcpaso{4}{milcMagenta}{\textbf{La ayuda concreta (8 min).} Al frente de cada fecha escribe \textbf{qué vas a hacer, no qué vas a decir}: \emph{acompañarla a algo, llevarle algo, invitarla a caminar, ayudarle con una materia, quedarte un rato en silencio}. \textbf{Y agrega una sola cosa dicha:} \emph{el nombre de la persona que murió, o de lo que se perdió, en una frase que no obligue a responder} ---«me acordé de él hoy»---.}

\begin{drawbox}{milcMagenta}{Antes de cerrar tu calendario}
\begin{checklista}
\item Tus dos columnas están llenas con hechos, no con impresiones.
\item Notaste si lo que ayudó fueron \textbf{acciones} y lo que no, \textbf{frases}.
\item Reemplazaste cuatro frases, y ningún reemplazo \textbf{le pide algo} a la persona.
\item Tu calendario tiene \textbf{tres momentos}, incluida una \textbf{fecha lejana}.
\item Cada fecha tiene \textbf{qué vas a hacer}, no qué vas a decir.
\item Entendiste que \textbf{nadie está atrasado} en su duelo.
\end{checklista}
\end{drawbox}

\begin{softbox}{milcMagenta}{milcGris}{Plantilla de trabajo}
\textbf{En vez de una frase hecha.} \emph{«No sé qué decir, pero aquí estoy.»} \\[1mm]
\textbf{En vez de «cualquier cosa me avisas».} \emph{«Paso el \_\_\_\_ a las \_\_\_\_.»} ---con día y hora---. \\[1mm]
\textbf{La que casi nadie usa.} \emph{«Hoy me acordé de \_\_\_\_.»} ---el nombre, sin pedir respuesta---.
\end{softbox}

\begin{infoband}{milcMagenta}


\textbf{En tu cuaderno (Actividad 3 · CREA).} Título: «Actividad 3. Mi calendario de acompañamiento». \textbf{Formato:} las dos columnas + las cuatro frases con su reemplazo + el calendario de tres momentos con la acción de cada uno. \textbf{Extensión:} una página. \textbf{Sabes que terminaste cuando} tienes \textbf{una fecha lejana} escrita en el calendario del celular.
\end{infoband}

\puente{El producto es un calendario. \emph{Porque «cualquier cosa me avisas» es una forma amable de no volver.}}

\milcestacion{milcMostaza}{Producto de la sesión}
\textbf{Calendario de acompañamiento: lo que ayudó, las frases reemplazadas y tres fechas con acciones concretas}.
\begin{softbox}{milcMostaza}{milcGris}{Criterios de calidad}
(1) La distinción entre lo que ayudó y lo que no se apoya en hechos concretos. (2) Los reemplazos de frases no exigen respuesta ni gratitud a quien está en duelo. (3) El calendario contiene tres momentos, uno de ellos lejano en el tiempo. (4) Cada fecha especifica una acción, no una frase. (5) Se explica, con evidencia, por qué no existen etapas que haya que cumplir.
\end{softbox}

\puente{Antes de cerrar, revisa lo que casi siempre falta: si tu plan no incluye una fecha lejana, se parece a todos los demás.}

\milcestacion{milcVino}{Evaluación liberadora · cinco dimensiones}

\begin{softbox}{milcVino}{milcGris}{Sentido de la evaluación}
Evaluar no es solo revisar una nota. Es reconocer qué aprendí, qué cambió en mí, qué emociones manejé, cómo me relaciono con la ciudad y la comunidad, y qué le dejaré a quien viene detrás.
\end{softbox}

\textbf{Mi reflexión liberadora en cinco dimensiones.} Completa cada frase iniciada:

\small
\begin{tabularx}{\textwidth}{>{\bfseries\RaggedRight\arraybackslash}p{3.4cm}Y>{\RaggedRight\arraybackslash}p{4.6cm}}
\rowcolor{milcVino}\color{white}Dimensión & \color{white}\bfseries Mi respuesta (frame starter) & \color{white}\bfseries Algo que descubrí\\
1. Personal & Lo que más cambió en mí fue \linead{6.5cm} & \linead{4.2cm}\\
2. Emocional & La emoción más fuerte fue \linead{2.5cm} y la regulé \linead{3cm} & \linead{4.2cm}\\
3. Ciudadana & En democracia esto importa porque \linead{6cm} & \linead{4.2cm}\\
4. Local (Cartago) & En mi barrio sería útil para \linead{6.5cm} & \linead{4.2cm}\\
5. Intergeneracional & A mi abuela/abuelo le diría \linead{6.5cm} & \linead{4.2cm}\\
\end{tabularx}
\normalsize

\begin{infoband}{milcVino}
\textbf{Para la próxima sustentación.} Vuelve a esta tabla en seis meses. Verás que las respuestas cambian — no porque lo aprendido cambie, sino porque tú cambias. La evaluación liberadora es una conversación contigo a través del tiempo.
\end{infoband}


\puente{Cerramos con tres voces: la comunidad que se funda escuchando, el cambio como ley de las cosas, y lo que las pantallas hicieron con la memoria.}

\milcestacion{milcGradeDark}{Triángulo de pensamiento · cierre filosófico}

\begin{softbox}{milcVino}{milcGris}{Dussel · lente del nosotros}
\textit{«Aquel que es capaz de escuchar silenciosamente al Otro como otro es el que puede constituir una comunidad y no una mera sociedad totalizada.»}

{\footnotesize\itshape\color{milcNegro!70}--- Enrique Dussel · Filosofía de la liberación (1977)}

\textbf{«Escuchar \emph{silenciosamente}».} \emph{En un duelo esa palabra deja de ser una metáfora y se vuelve una instrucción literal.} \textbf{Casi todo el daño que se hace acompañando viene de hablar:} \emph{la frase que consuela mal, el consejo, el lado bueno, la comparación.} \emph{Y casi todo lo que sirve viene del silencio acompañado ---estar ahí sin llenar el aire---.} \textbf{Fíjate también en «al Otro \emph{como otro}»:} \emph{cada quien hace su duelo distinto}. \textbf{Unos hablan sin parar y otros no quieren hablar; unos vuelven rápido a la rutina y otros no pueden; unos lloran y otros no lloran nunca ---y ninguno lo está haciendo mal---.} \emph{Quien acompaña de verdad no está esperando que el otro reaccione como él reaccionaría.}

\textbf{Cuando acompaño a alguien que está mal, ¿cuánto hablo yo?}
\end{softbox}
\linead{\linewidth}\\[1mm]
\linead{\linewidth}

\begin{softbox}{milcMostaza}{milcGris}{Estoicismo · Marco Aurelio · Meditaciones VII, 47 (c. 175 d.C.) · cuidado interior}
\textit{«Contempla las evoluciones de los astros y piensa al mismo tiempo que evolucionas con ellos; y reflexiona continuamente cómo los elementos cambian unos en otros.»}

{\footnotesize\itshape\color{milcNegro!70}--- Marco Aurelio · Meditaciones VII, 47 (c. 175 d.C.)}

\emph{Marco Aurelio mira lo que cambia para acordarse de que él también cambia.} \textbf{Y aquí hay que usarlo con mucho cuidado, porque tiene una versión mala:} \emph{decirle a alguien en duelo que «todo pasa» es apurarlo, y ya vimos que eso no ayuda}. \textbf{La versión buena es otra y va dirigida a quien está en medio de la pérdida:} \emph{cuando algo duele así, se siente permanente ---se siente que uno va a estar así siempre---}. \textbf{Y no es cierto.} \emph{No porque la ausencia desaparezca ---no desaparece--- sino porque \textbf{tú también vas a cambiar}, y la persona que serás en un año va a poder cargar esto de una manera que hoy no puedes.} \emph{Eso no se dice para consolar a otro. Se dice para uno mismo, en la noche difícil.}

\textbf{¿Estoy tratando esta pérdida como si el modo en que me siento hoy fuera para siempre?}
\end{softbox}
\linead{\linewidth}\\[1mm]
\linead{\linewidth}

\begin{softbox}{milcTurquesa}{milcGris}{Floridi · lente de la infoesfera}
\textit{«Las tecnologías digitales están re-ontologizando la naturaleza misma de la infoesfera.»}

{\footnotesize\itshape\color{milcNegro!70}--- Luciano Floridi · La cuarta revolución (2014; trad.)}

\textbf{El duelo cambió de forma, y esta generación es la primera que lo vive así.} \emph{Antes, de una persona que moría quedaban unas pocas fotos.} \textbf{Hoy queda un perfil que sigue ahí, con su voz en los audios, sus mensajes sin leer y un cumpleaños que la aplicación anuncia cada año.} \emph{Eso es re-ontologizar: cambió \textbf{lo que queda} de alguien.} \textbf{Y no es ni bueno ni malo por sí mismo ---depende de qué se haga con ello---:} \emph{para unos, releer los mensajes es una manera de seguir cerca; para otros, es quedarse en la puerta sin poder entrar ni salir.} \textbf{La pregunta útil no es si se debe borrar o guardar,} \emph{sino esta: cuando lo abro, ¿salgo mejor o peor de como entré?}

\textbf{Lo que guardo de lo que perdí, ¿me acompaña o me deja atrapado?}
\end{softbox}
\linead{\linewidth}\\[1mm]
\linead{\linewidth}

\begin{infoband}{milcNegro}
\textbf{Nota para el docente.} Las tres son citas textuales y llevan su referencia debajo. Puedes remitir a la obra en clase.
\end{infoband}

\milcestacion{milcVino}{Mi autoevaluación · marca una casilla por fila}

{\small
\setlength{\extrarowheight}{2.2pt}
\begin{tabularx}{\textwidth}{>{\RaggedRight\arraybackslash\bfseries}p{3.0cm}YYY}
\rowcolor{milcVino}\color{white}\bfseries Criterio & \color{white}\bfseries Lo logré & \color{white}\bfseries En proceso & \color{white}\bfseries Necesito apoyo\\[.6mm]
Comprensión del tema & \milcopcion{Explico las ideas con mis palabras.} & \milcopcion{Reconozco las ideas, falta precisión.} & \milcopcion{Necesito repasar lo básico.}\\[.8mm]
\rowcolor{milcGris}Calidad del acompañamiento & \milcopcion{Sé que no hay etapas que cumplir y tengo un plan de compañía con días concretos.} & \milcopcion{Entiendo que no se apura, pero solo se me ocurre decir cosas.} & \milcopcion{Sigo pensando que hay que ayudar a la gente a superarlo rápido.}\\[.8mm]
Aplicación y producto & \milcopcion{Mi producto cumple los criterios.} & \milcopcion{Está armado, con ajustes pendientes.} & \milcopcion{Aún está en construcción.}\\[.8mm]
\rowcolor{milcGris}Reflexión 5 dimensiones & \milcopcion{Respondí cada una con honestidad.} & \milcopcion{Respondí algunas con detalle.} & \milcopcion{Mis respuestas son muy generales.}\\[.8mm]
Triángulo de pensamiento & \milcopcion{Conecté las 3 voces con mi trabajo.} & \milcopcion{Conecté al menos una.} & \milcopcion{No las relaciono con mi proceso.}\\[.8mm]
\end{tabularx}}

\vspace{3mm}
\textbf{Hoy entendí que:}\\[1mm]
\linead{\linewidth}\\[1mm]
\linead{\linewidth}

\textbf{Mi compromiso para la próxima sesión es:}\\[1mm]
\linead{\linewidth}\\[1mm]
\linead{\linewidth}

\begin{softbox}{milcMostaza}{milcGris}{Cita de despedida · Marco Aurelio}
\textit{``El obstáculo en el camino se vuelve el camino. Lo que estorba al camino se vuelve camino.''}\\[1mm]
Cada error de hoy, bien observado, se convierte en pieza del camino que pisarás mañana.
\end{softbox}

\small
\begin{tabularx}{\textwidth}{>{\bfseries}p{3.6cm}Y}
\rowcolor{milcGradeDark!12}Nombre completo: & \linead{10cm}\\
Fecha: & \linead{4cm} \quad Curso/grupo: \linead{4cm}\\
Mi firma: & \linead{9cm}\\
\end{tabularx}
\normalsize

\begin{infoband}{milcGradeDark}
\textbf{Para la próxima sesión.} Dos cosas. \textbf{Primero, cumple la primera fecha de tu calendario} ---la de esta semana--- y \textbf{deja las otras dos puestas en el celular con alarma}. \emph{Anota qué hiciste y qué pasó; y si te atreviste a decir el nombre, anota cómo reaccionó.} \textbf{Segundo, pregunta en tu casa cómo se despedía la gente:} averigua con alguien mayor \textbf{qué se hacía cuando moría alguien} ---el velorio, el novenario, la novena, quién cocinaba, quién rezaba, cuántos días duraba, cuándo se consideraba terminado--- y \textbf{si hoy se sigue haciendo igual}. \textbf{Trae:} el resultado de tu primera fecha y lo que te contaron. \emph{Vas a encontrar casi siempre lo mismo, y es el argumento entero de esta guía: antes había un sistema ---días, horas, tareas repartidas, un final marcado--- y hoy queda muy poco de eso. No se perdió una costumbre bonita: se perdió la manera en que la gente se aseguraba de que nadie se quedara solo.}
\end{infoband}

\end{document}
